\SeventhArticle{13}
Um den Genauigkeitsgrad unserer Integralformel $RA+R'A'+R''A''+\text{usw.}+R^{(n)}A^{(n)}$ zu bestimmen, setzen wir allgemein
\[
Ra^m+R'a'{}^m+R''a''{}^m+\text{usw.}+R^{(n)}a^{(n)m}=\frac{1}{m+1}-k^{(m)}
\]
sodass $k^{(m)}$ die Differenz zwischen dem wahren und dem angenäherten Wert des Integrals $\int t^m\,\mathrm{d}t$ von $t=0$ bis $t=1$ ist. Entwickeln wir die einzelnen Brüche in Reihen, so erhalten wir\SourceLeaf{190}
\[
\begin{aligned}
&\frac{R}{t-a}+\frac{R'}{t-a'}+\frac{R''}{t-a''}+\text{usw.}+\frac{R^{(n)}}{t-a^{(n)}}\\
&\quad=(1-k)t^{-1}+(\tfrac12-k')t^{-2}+(\tfrac13-k'')t^{-3}\\ &\qquad+(\tfrac14-k''')t^{-4}+\text{usw.}\\
&\quad=t^{-1}+\tfrac12t^{-2}+\tfrac13t^{-3}+\tfrac14t^{-4}+\text{usw.}-\theta
\end{aligned}
\]
wenn wir setzen
\[
\theta=kt^{-1}+k't^{-2}+k''t^{-3}+k'''t^{-4}+\text{usw.}
\]
oder vielmehr (da wir bereits wissen, dass $k$, $k'$, $k''$, $k'''$ usw. bis $k^{(n)}$ von selbst verschwinden müssen)
\[
\theta=k^{(n+1)}t^{-(n+2)}+k^{(n+2)}t^{-(n+3)}+k^{(n+3)}t^{-(n+4)}+\text{usw.}
\]

Multiplikation mit $T$ ergibt
\[
\begin{aligned}&T\left(\frac{R}{t-a}+\frac{R'}{t-a'}+\frac{R''}{t-a''}+\text{usw.}+\frac{R^{(n)}}{t-a^{(n)}}\right)\\ &=T'+T''-T\theta\end{aligned}
\]
Die linke Seite dieser Gleichung ist ein Polynom in $t$ vom Grad $n$, und ihre Werte für $t=a$, $t=a'$, $t=a''$ usw. sind jeweils $MR$, $M'R'$, $M''R''$ usw. Da dasselbe für die Funktion $T'$ gilt, wie aus der Art der Bestimmung der Zahlen $R$, $R'$, $R''$ usw. selbst ersichtlich ist, muss die linke Seite der Gleichung mit $T'$ identisch sein, und daher gilt $T''=T\theta$. Somit entsteht $\theta$ durch Entwicklung des Bruchs $\frac{T''}{T}$; auf diese Weise lassen sich die Koeffizienten $k^{(n+1)}$, $k^{(n+2)}$ usw. beliebig weit bestimmen. Sind diese gefunden, so lautet die Korrektur unseres Näherungswertes des Integrals $\int y\,\mathrm{d}t$
\[
=k^{(n+1)}K^{(n+1)}+k^{(n+2)}K^{(n+2)}+\text{usw.}
\]
wenn die Reihe, in die $y$ entwickelt wird, lautet
\[
y=K+K't+K''tt+K'''t^3+\text{usw.}
\]

\SeventhArticle{14}
Zieht man es vor, die Korrektur durch die Koeffizienten der Reihe nach Potenzen von $t-\tfrac12$ auszudrücken,
\[
y=L+L'(t-\tfrac12)+L''(t-\tfrac12)^2+L'''(t-\tfrac12)^3+\text{usw.}
\]
so lautet sie\SourceLeaf{191}
\[
=l^{(n+1)}L^{(n+1)}+l^{(n+2)}L^{(n+2)}+l^{(n+3)}L^{(n+3)}+\text{usw.}
\]
wenn wir allgemein mit $l^{(m)}$ die Korrektur des Näherungswertes des Integrals $\int(t-\tfrac12)^m\,\mathrm{d}t$ bezeichnen. Diese Korrekturen $l^{(m)}$ sind mit den Korrekturen $k^{(m)}$ durch folgende Gleichung verknüpft:
\[
\begin{aligned}&l^{(m)}=k^{(m)}-\tfrac12mk^{(m-1)}+\tfrac14\cdot\frac{m\cdot m-1}{1\cdot2}k^{(m-2)}\\ &-\tfrac18\cdot\frac{m\cdot m-1\cdot m-2}{1\cdot2\cdot3}k^{(m-3)}+\text{usw.}\end{aligned}
\]
Um sie jedoch unabhängig davon zu bestimmen, beachten wir: Die Funktion $\theta$ wird durch die Substitution $t=\tfrac12u+\tfrac12$ zu
\[
\begin{aligned}
&\phantom{{}+{}}2k(u^{-1}-u^{-2}+u^{-3}-u^{-4}+\text{usw.})\\
&+4k'(u^{-2}-2u^{-3}+3u^{-4}-4u^{-5}+\text{usw.})\\
&+8k''(u^{-3}-3u^{-4}+6u^{-5}-10u^{-6}+\text{usw.})\\
&+16k'''(u^{-4}-4u^{-5}+10u^{-6}-20u^{-7}+\text{usw.})\\
&+\text{usw.}
\end{aligned}
\]
oder zu
\[
\begin{aligned}
&2ku^{-1}+4(k'-\tfrac12)u^{-2}+8(k''-\tfrac12\cdot2k'+\tfrac14k)u^{-3}\\
&\qquad+16(k'''-\tfrac12\cdot3k''+\tfrac14\cdot3k'-\tfrac18k)u^{-4}+\text{usw.}
\end{aligned}
\]
oder zu
\[
2lu^{-1}+4l'u^{-2}+8l''u^{-3}+16l'''u^{-4}+\text{usw.}
\]
oder schließlich, da wir im Voraus wissen, dass $l$, $l'$, $l''$, $l'''$ usw. bis $l^{(n)}$ von selbst verschwinden, zu
\[
\begin{aligned}&2^{n+2}l^{(n+1)}u^{-(n+2)}+2^{n+3}l^{(n+2)}u^{-(n+3)}\\ &+2^{n+4}l^{(n+4)}u^{-(n+4)}+\text{usw.}\end{aligned}
\]
Es ist aber $\theta=\frac{T''}{T}$. Da $T$, $T''$ durch die Substitution $t=\tfrac12u+\tfrac12$ in $\frac{U}{2^{n+1}}$, $\frac{U''}{2^n}$ übergehen (Artikel 9), geht die Funktion $\theta$ durch dieselbe Substitution in $\frac{2U''}{U}$ über. Bezeichnen wir daher die durch Entwicklung des Bruchs $\frac{U''}{U}$ entstehende Reihe mit $\Omega$, so ist
\[
\begin{aligned}&\Omega=2^{n+1}l^{(n+1)}u^{-(n+2)}+2^{n+2}l^{(n+2)}u^{-(n+3)}\\ &+2^{n+3}l^{(n+3)}u^{-(n+4)}+\text{usw.}\end{aligned}
\]
wodurch sich die Koeffizienten $l^{(n+1)}$, $l^{(n+2)}$ usw. beliebig weit bestimmen lassen.

So finden wir im Beispiel des Artikels 10
\[
\begin{aligned}
U''&=-\tfrac{176}{13125}u^{-1}-\tfrac{304}{28125}u^{-3}-\tfrac{2576}{309375}u^{-5}-\text{usw.}\\
\Omega&=-\tfrac{176}{13125}u^{-7}-\tfrac{832}{28125}u^{-9}-\tfrac{189856}{4296875}u^{-11}-\text{usw.}
\end{aligned}
\]
\SourceLeaf{192}und daher lautet die Korrektur des Näherungswertes des Integrals
\[
=-\frac{11}{52500}L^{\mathrm{VI}}-\frac{13}{112500}L^{\mathrm{VIII}}-\frac{5933}{137500000}L^{\mathrm{X}}-\text{usw.}
\]
\SeventhArticle{15}
Der Koeffizient $K^{(m)}$ der in eine Reihe entwickelten Funktion $y$ ist nach dem Satz von \textsc{Taylor} gleich dem Wert von
\[
\frac{1}{1.2.3\mathinner{.\,.\,.\,.}m}\mathbin{.}\frac{\mathrm{d}^m y}{\mathrm{d}t^m}
\quad\text{oder}\quad
\frac{\Delta^m}{1.2.3\mathinner{.\,.\,.\,.}m}\mathbin{.}\frac{\mathrm{d}^m y}{\mathrm{d}x^m}
\]
für $t=0$ oder $x=g$; entsprechend ist der Koeffizient $L^{(m)}$ der Wert desselben Ausdrucks für $t=\frac12$ oder $u=0$ oder $x=g+\frac12\Delta$. Beiden Koeffizienten schreiben wir die \textit{Ordnung} $m$ zu. Im Allgemeinen ist unsere Integration daher bis zur Ordnung $n$ einschließlich exakt, welche Werte auch für $a,a',a''\mathinner{.\,.\,.\,.}a^{(n)}$ genommen werden. Dies hindert jedoch nicht daran, durch geschickte Wahl dieser Werte die Genauigkeit auf einen höheren Grad zu bringen. So haben wir bereits oben gesehen, dass sich beim Verfahren von \textsc{Cotes}, also für $a=0,a'=\frac1n,a''=\frac2n,a'''=\frac3n$ usw., die Genauigkeit von selbst bis zur Ordnung $n+1$ einschließlich erstreckt, wenn $n$ eine gerade Zahl ist. Allgemein ist klar: Werden die Werte $a,a',a'',a'''$ usw. so gewählt, dass in der Funktion $T''$ oder $U''$ ein oder mehrere Anfangsglieder wegfallen, so wird die Genauigkeit um ebenso viele Ordnungen über die Ordnung $n$ hinaus erhöht, wie Glieder wegfallen. Da die Anzahl der Größen, die wir frei wählen dürfen, $n+1$ beträgt, folgt daraus leicht, dass sich durch ihre geeignete Bestimmung die Genauigkeit stets bis zur Ordnung $2n+1$ einschließlich steigern lässt. Auf diese Weise lässt sich mit $n+1$ Gliedern dieselbe Genauigkeitsordnung erreichen, für die sonst $2n+1$ oder $2n+2$ Glieder verwendet werden müssten, wenn wir dem Verfahren von \textsc{Cotes} folgten.

